\section{Spatial transfer with endpoint fields}\label{sec:transfer}
We first put the periodic and open chains in the same transfer
representation. The boundary vectors, including their complex phases,
will distinguish the physical endpoint fields.
The time-ordered bond expansion has a general formulation in
\cite[arXiv v1, Section 2.2, Equations (2.10)--(2.12)]{AizenmanNachtergaele1994}.
We prove the required boundary identity directly with the specified
kernel and field coefficients.

\subsection{Hamiltonians and shifted partition functions}
From now on set
\begin{equation}\label{tr:constants}
 h=\frac35,\qquad a=\frac{700741}{500000},\qquad
 b_0=\frac{49}{4},\qquad b_j=2^jb_0\quad(j\ge0).
\end{equation}
It is convenient to label an $n$-site chain by $0,\ldots,n-1$.
Write $S_j\cdot S_k=\sum_\alpha S_j^\alpha S_k^\alpha$ and define
\begin{align}
 H_n^{\rm opp}&=\sum_{j=0}^{n-2}S_j\cdot S_{j+1}
                         -hS_0^z+hS_{n-1}^z,\label{tr:Hopp}\\
 H_n^{\rm s}&=\sum_{j=0}^{n-2}S_j\cdot S_{j+1}
                         -hS_0^z-hS_{n-1}^z,\label{tr:Hsame}\\
 H_n^{\rm per}&=\sum_{j=0}^{n-2}S_j\cdot S_{j+1}
                         +S_{n-1}\cdot S_0\quad(n\ge4).
\end{align}
For $n=2L+1$, $H_n^{\rm s}$ is \eqref{eq:target} after relabeling.
Our partition functions are ordinary traces over the entire physical
Hilbert space:
\begin{equation}\label{tr:partitions}
 \begin{split}
 Z_n(b)&=\Tr e^{-b(H_n^{\rm per}+an\Id)},\\
 Y_n(b)&=\Tr e^{-b(H_n^{\rm opp}+an\Id)},\qquad
 Y_n^{\rm s}(b)=\Tr e^{-b(H_n^{\rm s}+an\Id)}.
 \end{split}
\end{equation}
The shift is per site, also for open chains. It will cancel from all
physical Gibbs weight ratios.

\subsection{The spatial operator}
We use the concrete spatial-transfer operator in
\cite[Proposition 3.1]{periodic}. Its auxiliary space records ordered
bond events in imaginary time; each physical site contributes a kernel
depending on its two incident histories. Use the unitarily equivalent
Cartesian
spin matrices $(S^\alpha)_{rs}=-i\varepsilon_{\alpha rs}$ and set
$G_\alpha=iS^\alpha$. Then $G_\alpha$ is real,
$\norm{G_\alpha}=1$, and
\[
 -S\cdot S=\sum_{\alpha=x,y,z}G_\alpha\otimes G_\alpha.
\]
A history is an ordered list
$\omega=((t_1,\alpha_1),\ldots,(t_d,\alpha_d))$ with
$0<t_1<\cdots<t_d<b$ and $\alpha_r\in\{x,y,z\}$; the empty list
is allowed. Give the disjoint union $\Omega_b$ of these labeled
simplexes Lebesgue measure on each simplex and mass one at the empty
history. Its measure $\mu_b$ has total mass $e^{3b}$.

For two histories, merge their events in chronological order and take
the trace of the resulting product of $G$ matrices, with increasing
time from left to right. Denote this real kernel by $k_b(\omega,\eta)$.
Set it to zero at nonempty time
collisions, a null set, and use trace $3$ for the two empty histories.
The kernel is symmetric: exchanging the two histories leaves their
merged list unchanged. In particular, symmetry does not reverse a
product of noncommuting matrices. Moreover $|k_b|\le3$.
On $\mathcal K_b=L^2(\Omega_b,\mu_b;\C)$ define
\begin{equation}\label{tr:operator}
 (X_bf)(\omega)=e^{-ab}\int_{\Omega_b}k_b(\omega,\eta)
                                      f(\eta)\,d\mu_b(\eta).
\end{equation}
It is real, self-adjoint, and Hilbert--Schmidt. The cited transfer
Proposition identifies its periodic moments:
\begin{equation}\label{tr:periodic}
 Z_n(b)=\Tr X_b^n\quad(n\ge4),\qquad
 \sum_i|\mu_i(b)|^k<\infty\quad(k\ge2),
\end{equation}
where $\mu_i(b)$ lists its nonzero eigenvalues with multiplicity.
This exact operator, with the shift in \eqref{tr:constants}, is the one
used for the periodic inputs in Section~\ref{sec:periodic}.

Rotations will let us change one endpoint field and, later, show that
the leading spatial line is fixed. Let $D_2$ consist of the identity
and the three diagonal rotations by
$\pi$, let $C$ cyclically permute the coordinate axes, and put
$\mathcal G=\langle D_2,C\rangle$. This is a group of twelve proper
signed permutations. If $ge_\alpha=\epsilon_g(\alpha)
e_{\sigma_g(\alpha)}$, write $\sigma_g\omega$ for relabeling the
marks and $s_g(\omega)=\prod_r\epsilon_g(\alpha_r)$. The operators
\begin{equation}\label{tr:rotations}
 (U_gf)(\omega)=s_g(\sigma_g^{-1}\omega)f(\sigma_g^{-1}\omega)
\end{equation}
form a real unitary representation and commute with $X_b$. This follows
either directly by rotating every matrix in the kernel, or from the
same Proposition. In particular, fix
$P=\diag(1,-1,-1)$, so $Pe_z=-e_z$.

\subsection{Boundary vectors}
Define $v_b\in\mathcal K_b$ by
\begin{equation}\label{tr:vector}
 v_b(\omega)=
 \begin{cases}
 (ih)^d,&\text{all $d$ marks of $\omega$ equal $z$},\\
 0,&\text{otherwise}.
 \end{cases}
\end{equation}
The empty-history value is one. The simplex volumes give
$\norm{v_b}^2=\sum_{d\ge0}b^dh^{2d}/d!=e^{bh^2}$.
The inner product below is conjugate linear in its first argument.

\begin{proposition}[Open-chain transfer]\label{tr:boundary}
For $b>0$, $n\ge2$, and $g\in\mathcal G$, the pairing
$\langle v_b,X_b^nU_gv_b\rangle$ is the shifted partition function
of the open chain with left field $-hS_0^z$ and right field
$+hS_{n-1}^{ge_z}$. Consequently
\begin{equation}\label{tr:open}
 Y_n(b)=\langle v_b,X_b^nv_b\rangle,\qquad
 Y_n^{\rm s}(b)=\langle v_b,X_b^nU_Pv_b\rangle.
\end{equation}
Here $S^u=\sum_\alpha u_\alpha S^\alpha$ for a real vector $u$.
\end{proposition}
\begin{proof}
Expand the negative Hamiltonian exponential into its ordered-time
series. Each bond event inserts two $G_\alpha$ matrices, by the
Cartesian identity above. At the left endpoint, the conjugated
coefficient in \eqref{tr:vector} gives
$(-ih)G_z=+hS^z$ per event; at the right endpoint the coefficient is
$(ih)G_z=-hS^z$. Relabeling by $g$ and multiplying its signs changes
this last insertion to $-hS^{ge_z}$. These are exactly the two field
insertions in the negative Hamiltonian stated in the Proposition.

Regroup the series by the boundary histories $\omega_0,\omega_n$
and the $n-1$ bond histories $\omega_1,\ldots,\omega_{n-1}$.
Their interleavings partition each full time simplex up to null
collisions. At site $r$, $0\le r<n$, the trace is
$k_b(\omega_r,\omega_{r+1})$, in the original chronological order.
Integrating these $n$ kernel factors against the two boundary
coefficients gives $X_b^n$ and the shift factor $e^{-ban}$.
No bond joins the endpoints.

This regrouping is legitimate: the absolute sum of traced terms is
bounded by
$3^n\exp\{b[3(n-1)+2h]\}$ before the scalar shift. All sums and
integrals are therefore absolutely convergent. Choosing $g=1$ and
$g=P$ gives \eqref{tr:open}.
\end{proof}

\begin{figure}[!htbp]
\centering
\begin{tikzpicture}[>=Stealth,
  site/.style={draw,minimum width=9mm,minimum height=8mm},
  every node/.style={font=\small}]
 \foreach \y in {0,-1.8} {
   \node[site] at (0,\y) {$X_b$};
   \node[site] at (1.6,\y) {$X_b$};
   \node at (3.2,\y) {$\cdots$};
   \node[site] at (4.8,\y) {$X_b$};
   \draw (-1.6,\y)--(-.45,\y);
   \draw (.45,\y)--(1.15,\y);
   \draw (2.05,\y)--(2.75,\y);
   \draw (3.65,\y)--(4.35,\y);
   \draw (5.25,\y)--(6.4,\y);
   \node[left] at (-1.6,\y) {$\overline v_b$};
 }
 \node[right] at (6.4,0) {$v_b$};
 \node[right] at (6.4,-1.8) {$U_Pv_b$};
 \node[above] at (-.2,.45) {left: $-hS^z$};
 \node[above] at (5.0,.45) {right: $+hS^z$};
 \node[below] at (-.2,-2.25) {left: $-hS^z$};
 \node[below] at (5.0,-2.25) {right: $-hS^z$};
 \draw[->] (6.7,-.45)--(6.7,-1.35);
 \node[left] at (6.45,-.9) {$P e_z=-e_z$};
 \node at (2.4,-.65) {$n$ site kernels};
\end{tikzpicture}
\caption{The two spatial boundary pairings. Each box contributes one
site kernel; the internal integrations represent the $n-1$ bonds.
Replacing the right boundary vector by $U_Pv_b$ changes the right
physical field. This is an identity between transfer representations,
not a claim that the two physical Hamiltonians are unitarily equivalent.}
\label{fig:boundary-pairings}
\end{figure}

For even lengths, the opposite-field expression is the positive moment
$Y_n(b)=\langle v_b,|X_b|^n v_b\rangle$.
At odd lengths, the powers of $X_b$ retain the signs of its eigenvalues.
We first concentrate the even opposite-field moments. The resulting
leading-line estimate will then control the signed same-field pairing
at odd lengths; Figure~\ref{fig:boundary-pairings} distinguishes the
two boundary pairings.
